\origpage{253}{GRAPH MINORS. XII}
\hspace*{12pt}(1) \textit{Let $\Delta_1, \dots, \Delta_t \subseteq \Sigma$ be disjoint closed discs, let $N_i \subseteq bd(\Delta_i)$ with $|N_i| = 4$ for $1 \le i \le t$, and let $Z = N_1 \cup \dots \cup N_t$. Then there is a partition $(Z_v \colon v \in V(L))$ of $Z$ satisfying the topological feasibility condition of Section 3, such that for $1 \le i \le p$, if $e, e' \in E(L)$ cross at $x_i$ and have ends $u, v$ and $u', v'$, respectively in $L$, then the four members of $N_i$ can be enumerated in order on $bd(\Delta_i)$ as $\{a, b, c, d\}$, where $a \in Z_u, b \in Z_{u'}, c \in Z_v, d \in Z_{v'}$.}

Let us choose $\theta$ to satisfy (3.2) for $\Sigma$ and $t$, taking $z = 4t$. We claim that $\theta$ satisfies (4.5). For let $G, \Gamma, \mathscr{T}, r_i, C_i, a_i, b_i, c_i, d_i, P_i, Q_i$ ($1 \le i \le t$) be as in (4.5). For $1 \le i \le t$ let $N_i = \{a_i, b_i, c_i, d_i\}$, and let $Z = N_1 \cup \dots \cup N_t$. Let $(Z_v \colon v \in V(L))$ be as in (1). By (3.2), there are disjoint connected subgraphs $H_v$ ($v \in V(L)$) of $\Gamma$ such that $V(H_v) \cap Z = Z_v$ ($v \in V(L)$).

Let us renumber $P_1, Q_1, \dots, P_t, Q_t$ as $R_e$ ($e \in E(L)$), in such a way that if $e$ is involved in $x_i$ and has ends $u, v$, then $R_e$ is one of $P_i, Q_i$ and has ends in $V(H_u), V(H_v)$, respectively. (This is possible since $(Z_v \colon v \in V(L))$ satisfies (1).) For each $e \in E(L)$ let $\alpha(e) \in E(R_e)$. For each $v \in V(L)$ let $\alpha(v)$ be the component of
\begin{equation*}
(\bigcup (H_v \colon v \in V(L)) \cup \bigcup (R_e \colon e \in E(L))) \setminus \{\alpha(e) \colon e \in E(L)\}
\end{equation*}
which includes $H_v$. Then it follows that if $e \in E(L)$ has ends $u, v$ then $\alpha(e)$ has ends in $V(\alpha(u)), V(\alpha(v))$.

Consequently, conditions (i)--(iv) in the definition of ``controls an $L$ minor'' are satisfied. For (v), let $\mathscr{T}^*$ be the tangle in $G$ induced by $\mathscr{T}$. Then $\mathscr{T}^*$ has order $\ge \theta \ge 2t \ge |V(L)|$; let $(A, B) \in \mathscr{T}^*$ and suppose that $V(\alpha(v)) \subseteq V(A)$ for some $v \in V(L)$. There are two edges of $L$ incident with $v$ that do not cross each other, and so there exist $i, j$ with $1 \le i < j \le t$ such that $Z_v \cap N_i, Z_v \cap N_j \ne \emptyset$. Choose $a_i \in Z_v \cap N_i, a_j \in Z_v \cap N_j$. Since $d(r_i, r_j) \ge \theta$ it follows that $d(a_i, a_j) \ge \theta - 2$, and so (applying (4.1) as in the proof of (4.3) and (4.4)) $(A \cap \Gamma, B \cap \Gamma) \in \mathscr{T}$ has order $\ge \frac{1}{2}(\theta - 2) \ge |V(L)|$. Consequently, $\mathscr{T}^*$ controls an $L$ minor of $G$, and since $L$ has a $K_p$ minor, it follows from (4.2) that $\mathscr{T}^*$ controls a $K_p$ minor of $G$, as required. $\square$

\sectionhead{5. \textsc{Preserving} \textsc{Respect}}

The remainder of the paper is concerned with the effect on a respectful tangle and the associated metric of making local changes in the drawing or surface. The first local change we need to study is the following. Suppose we have a drawing in a surface, with a respectful tangle. Now we add a ``handle'' to the surface, attaching it on at two regions of the drawing, and we extend the drawing onto the handle. We know that the original tangle induces a tangle in the enlarged drawing; our question here is, under what
